\documentclass{article}
\title{Resumos de Cálculo Diferencial e Integral III}
\author{Gabriel Luca Silva}
\date{}

\usepackage{amsmath}
\usepackage{amssymb}
% \usepackage{bbm}
% \usepackage{comment}
\usepackage{geometry}
\usepackage{graphicx}
\usepackage{hyperref}
\hypersetup{
    colorlinks=false,
    pdfborder={0 0 0}
}
\usepackage{indentfirst}
\usepackage[utf8]{inputenc}
\usepackage{mathtools}
\usepackage{subcaption}

\geometry{
    a4paper,
    % total={129mm,170mm},
    left=23mm,
    top=23mm,
    bottom=28mm,
    right=28mm,
}

% -------------------------------------------------------

\begin{document}
\maketitle
\vspace*{\fill}
\textbf{Estes resumos não são totalmente rigorosos e destinam-se principalmente a servir como
um guia de consulta rápida, em vez de um texto explicativo detalhado.}
\bigskip


\newpage


\renewcommand{\contentsname}{Índice}
\tableofcontents


\newpage


\section{Introdução}
O foco principal de Cálculo Diferencial e Integral III são as equações diferenciais.
Estes resumos apresentam métodos de resolução para os seus diferentes tipos.
Faz sentido começar por algumas definições.

\subsection{Equação Diferencial}
Uma equação diferencial é uma equação que envolve uma função desconhecida e as suas derivadas.

As incógnitas destas equações são funções, por isso
as soluções das equações diferenciais vão ser sempre funções.

\subsection{Equação Diferencial Ordinária (EDO)}
A função desconhecida depende de uma única variável independente.

$$
\frac{d^2 x}{dt^2} + \omega^2 x = 0 \quad \text{que corresponde a} \quad x'' + \omega^2 x = 0
$$

\begin{itemize}
    \item A função desconhecida é $x(t)$.
    \item A função $x(t)$ é uma variável dependente que depende apenas da variável independente $t$.
\end{itemize}

\subsection{Equação diferencial parcial (EDP)}
A função desconhecida depende de várias variáveis independentes.

$$
\frac{\partial u}{\partial t} = \alpha \frac{\partial^2 u}{\partial x^2}
$$

\begin{itemize}
    \item A função desconhecida é $u(x,t)$.
    \item A função $u(x,t)$ é uma variável dependente que depende das variáveis independentes $x$ e $t$.
\end{itemize}

\subsection{Ordem de uma Equação Diferencial}
A ordem de uma equação diferencial é a maior ordem de derivada presente na equação.

$$
y''' + 3y'' - 4y' + 7y = 0
$$

\begin{itemize}
    \item A equação diferencial é de ordem 3.
\end{itemize}

\subsection{Equação Diferencial Ordinária Linear e Não Linear}
Uma equação diferencial ordinária de ordem $n$ é dita linear se pode ser escrita na forma:

$$
a_0(t) y^{(n)} + a_1(t) y^{(n-1)} + \ldots + a_{n-1}(t) y' + a_n(t) y = g(t)
$$

onde $a_n(t), a_{n-1}(t), \ldots, a_1(t), a_0(t)$ e $g(t)$ são funções conhecidas.\par
\bigskip
Se a equação não pode ser escrita nesta forma, é dita não linear.

\subsubsection{Equação Diferencial Ordinária Linear Homogénea e Não Homogénea}
Uma equação diferencial ordinária linear de ordem $n$ pode ser classificada como:

\begin{itemize}
    \item Homogénea se $g(t) = 0$.
    \item Não Homogénea se $g(t) \neq 0$.
\end{itemize}


\newpage


\section{Equações Diferenciais Ordinárias de Primeira Ordem}

\subsection{Equação Diferencial Ordinária de 1ª Ordem Linear Homogénea}
\label{subsec: edo-ordem1-linear-homogenea}

\begin{align*}
& \boxed{y' + a(t) y = 0} \\
& y' = -a(t) y \\
& \frac{y'}{y} = -a(t) \\
& \frac{d}{dt} \Big[ \ln|y| \Big] = -a(t) \\
& \ln|y| = \int -a(t) dt \\
& \ln|y| = - \int a(t) dt + c, \quad c \in \mathbb{R} \\
& |y| = e^{-\int a(t) dt + c}, \quad c \in \mathbb{R} \\
& |y| = e^c e^{-\int a(t) dt}, \quad c \in \mathbb{R} \\
& |y| = k e^{-\int a(t) dt}, \quad k \in \mathbb{R}^+ \text{ pois } k = e^c > 0 \\
& y = \pm k e^{-\int a(t) dt}, \quad k \in \mathbb{R}^+ \\
& y = k e^{-\int a(t) dt}, \quad k \in \mathbb{R} \setminus \{0\} \\
& \text{Como } y(t) = 0 \text{ é também solução, conclui-se que:} \\
& \boxed{y(t) = k e^{-\int a(t) dt}, \quad k \in \mathbb{R}}
\end{align*}

\bigskip

\subsection{Equação Diferencial Ordinária de 1ª Ordem Linear Não Homogénea}
\label{subsec: edo-ordem1-linear-nao-homogenea}

Para este tipo de equação, usamos o \textbf{método do fator integrante}.
Este método introduz uma função auxiliar $\mu(t)$, a que se chama fator integrante,
com o objetivo de transformar o lado esquerdo da equação numa derivada do produto de duas funções, $\mu(t)$ e $y(t)$.

\begin{gather*}
\boxed{y' + a(t) y = g(t)} \\
\mu(t) y' + \mu(t) a(t) y = \mu(t) g(t)
\end{gather*}

Queremos escrever o lado esquerdo como

\begin{gather*}
\frac{d}{dt} \Big[ \mu y \Big] = \mu' y + \mu y'
\end{gather*}

Para isso, é necessário que

\begin{gather*}
\mu' = \mu a(t)
\end{gather*}

Sendo esta equação uma \nameref{subsec: edo-ordem1-linear-homogenea}, tal que

\begin{gather*}
\mu' + (-a(t)) \mu = 0 \\
\implies \mu(t) = e^{\int a(t) dt}
\end{gather*}

Então, calculado o fator integrante, podemos reescrever a equação original como

\begin{align*}
& \frac{d}{dt} \Big[ \mu(t) y \Big] = \mu(t) g(t) \\
& \mu(t) y= \int \mu(t) g(t) dt \\
& \mu(t) y = \int \mu(t) g(t) dt + c, \quad c \in \mathbb{R} \\
& y(t) = \frac{1}{\mu(t)} \left( \int \mu(t) g(t) dt + c \right), \quad c \in \mathbb{R} \\
& \boxed{y(t) = e^{-\int a(t) dt} \left( \int e^{\int a(t) dt} g(t) dt + c \right), \quad c \in \mathbb{R}}
\end{align*}

\bigskip

\subsection{Equação Diferencial Ordinária de 1ª Ordem Separável}

Uma equação diferencial é dita separável se pode ser escrita na forma

\begin{gather*}
\boxed{\frac{dy}{dt} = \frac{g(t)}{f(y)}}
\end{gather*}

Separando as variáveis, obtemos

\begin{gather*}
f(y) \frac{dy}{dt} = g(t)
\end{gather*}

Integrando ambos os lados em ordem a $t$, obtemos

\begin{gather*}
\int f(y) \frac{dy}{dt} dt = \int g(t) dt
\end{gather*}

Ficando com $y(t)$ definido implicitamente por

\begin{gather*}
\boxed{\int f(y) dy = \int g(t) dt}
\end{gather*}

\smallskip

\begin{itemize}
    \item Uma Equação Separável pode ser linear ou não linear.
    \item Uma \nameref{subsec: edo-ordem1-linear-homogenea} é separável.
    \item Uma \nameref{subsec: edo-ordem1-linear-nao-homogenea} não é separável.
\end{itemize}

\bigskip

\subsection{Equação Diferencial Ordinária de 1ª Ordem Exata}
\label{subsec: edo-ordem1-exata}

\begin{gather*}
\boxed{M(t,y) + N(t,y) \frac{dy}{dt} = 0}
\end{gather*}

Esta equação diferencial é dita exata se puder ser escrita como

\begin{gather*}
\frac{d}{dt} \Big[ \phi(t,y) \Big] = 0
\end{gather*}

Usando a regra da cadeia, temos que

\begin{gather*}
\frac{d}{dt} \Big[ \phi(t,y) \Big] = \frac{\partial \phi}{\partial t} + \frac{\partial \phi}{\partial y} \frac{dy}{dt}
\end{gather*}

Logo, para que a equação seja exata, deve ser possível encontrar uma função $\phi(t,y)$ tal que

\begin{gather*}
\frac{\partial \phi}{\partial t} = M(t,y) \quad \text{ e } \quad \frac{\partial \phi}{\partial y} = N(t,y)
\end{gather*}

A solução da equação exata é dada implicitamente por

\begin{gather*}
\boxed{\phi(t,y) = c, \quad c \in \mathbb{R}}
\end{gather*}

\begin{itemize}
    \item Nota: uma equação exata pode ser linear ou não linear.
\end{itemize}

\smallskip

\paragraph{Teorema}
\begin{gather*}
\text{Existe } \phi(t,y) \text{ tal que } \frac{\partial \phi}{\partial t} = M(t,y) \text{ e } \frac{\partial \phi}{\partial y} = N(t,y) \text{ se, e só se, } \frac{\partial M}{\partial y} = \frac{\partial N}{\partial t}
\end{gather*}


\paragraph{Método de resolução de uma equação exata}
\label{par: metodo-resolucao-edo-exata}
\begin{gather*}
M(t,y) + N(t,y) \frac{dy}{dt} = 0
\end{gather*}

\begin{enumerate}
    \item Verificar se é exata.
    \begin{gather*}
    \frac{\partial M}{\partial y} = \frac{\partial N}{\partial t}
    \end{gather*}
    
    \item Se for exata, encontrar a função $\phi(t,y)$, tal que
    \begin{gather*}
    \begin{dcases}
    \frac{\partial \phi}{\partial t} = M \implies \phi = \int M dt + c_1(y) \\
    \frac{\partial \phi}{\partial y} = N \implies \phi = \int N dy + c_2(t)
    \end{dcases}
    \end{gather*}
    
    Basta calcular uma das equações, e depois determinar o seu termo $c$.
    
    Por exemplo, podemos calcular a primeira equação,
    \begin{gather*}
    \phi = \int M dt + c_1(y)
    \end{gather*}

    E depois para determinar $c_1(y)$, usamos a segunda equação com o $\phi$ encontrado.
    \begin{gather*}
    \frac{\partial \phi}{\partial y} = N \\
    \frac{\partial }{\partial y} \left( \int M dt + c_1(y) \right) = N \\
    c_1'(y) = N - \frac{\partial }{\partial y} \left( \int M dt \right)
    \end{gather*}

    \item A solução é dada implicitamente por
    \begin{gather*}
    \phi(t,y) = c, \quad c \in \mathbb{R}
    \end{gather*}
\end{enumerate}

\bigskip

\subsection{Equação Diferencial Ordinária de 1ª Ordem Redutível a Exata}

Por vezes, uma equação diferencial que não é exata pode ser transformada numa equação exata 
usando o \textbf{método do fator integrante} (visto anteriormente na resolução de uma \nameref{subsec: edo-ordem1-linear-nao-homogenea})

\begin{gather*}
\boxed{M(t,y) + N(t,y) \frac{dy}{dt} = 0}
\end{gather*}

Se a equação diferencial não for exata, ou seja, se

\begin{gather*}
\frac{\partial M}{\partial y} \neq \frac{\partial N}{\partial t}
\end{gather*}

Então vamos procurar um fator integrante $\mu(t,y)$ para tornar a equação exata.

\begin{gather*}
\mu(t,y) M(t,y) + \mu(t,y) N(t,y) \frac{dy}{dt} = 0
\end{gather*}

Para que esta nova equação seja exata, é necessário que

\begin{gather*}
\frac{\partial }{\partial y} \Big[ \mu M \Big] = \frac{\partial }{\partial t} \Big[ \mu N \Big] \\
\frac{\partial \mu}{\partial y} M + \mu \frac{\partial M}{\partial y} = \frac{\partial \mu}{\partial t} N + \mu \frac{\partial N}{\partial t}
\end{gather*}

Esta é uma equação difícil de resolver para determinar o fator integrante $\mu(t,y)$.
No entanto, em alguns casos, podemos encontrar um fator integrante que depende apenas de $t$ ou apenas de $y$.

\bigskip

Para $\mu = \mu(t)$, a equação torna-se

\begin{gather*}
0 + \mu \frac{\partial M}{\partial y} = \mu' N + \mu \frac{\partial N}{\partial t} \\
N \mu' = \mu \left( \frac{\partial M}{\partial y} - \frac{\partial N}{\partial t} \right) \\
\mu' = \mu \frac{\frac{\partial M}{\partial y} - \frac{\partial N}{\partial t}}{N} \\
\end{gather*}

Esta é uma \nameref{subsec: edo-ordem1-linear-homogenea}, que pode ser resolvida se, e só se

\begin{gather*}
\frac{\frac{\partial M}{\partial y} - \frac{\partial N}{\partial t}}{N} \quad \text{depender apenas de } t
\end{gather*}

Se assim for, ficamos então com

\begin{gather*}
\mu' + \left( - \frac{\frac{\partial M}{\partial y} - \frac{\partial N}{\partial t}}{N} \right) \mu = 0 \\
\implies \mu(t) = e^{\int \frac{\frac{\partial M}{\partial y} - \frac{\partial N}{\partial t}}{N} dt}
\end{gather*}

Similarmente, para $\mu = \mu(y)$, a equação torna-se
\begin{gather*}
\mu' M + \mu \frac{\partial M}{\partial y} = 0 + \mu \frac{\partial N}{\partial t} \\
M \mu' = \mu \left( \frac{\partial N}{\partial t} - \frac{\partial M}{\partial y} \right) \\
\mu' = \mu \frac{\frac{\partial N}{\partial t} - \frac{\partial M}{\partial y}}{M} \\
\end{gather*}

Esta é uma \nameref{subsec: edo-ordem1-linear-homogenea}, que pode ser resolvida se, e só se

\begin{gather*}
\frac{\frac{\partial N}{\partial t} - \frac{\partial M}{\partial y}}{M} \quad \text{depender apenas de } y
\end{gather*}

Se assim for, ficamos então com

\begin{gather*}
\mu' + \left( - \frac{\frac{\partial N}{\partial t} - \frac{\partial M}{\partial y}}{M} \right) \mu = 0 \\
\implies \mu(y) = e^{\int \frac{\frac{\partial N}{\partial t} - \frac{\partial M}{\partial y}}{M} dy}
\end{gather*}

\bigskip

\begin{itemize}
    \item Se conseguirmos encontrar um fator integrante $\mu(t)$ ou $\mu(y)$, fica-se então com uma \nameref{subsec: edo-ordem1-exata} que pode ser resolvida como visto anteriormente.
\end{itemize}

\paragraph{Método de resolução de uma equação redutível a exata}
\begin{gather*}
M(t,y) + N(t,y) \frac{dy}{dt} = 0
\end{gather*}

\begin{enumerate}
    \item Verificar se é exata.
    \begin{gather*}
    \frac{\partial M}{\partial y} = \frac{\partial N}{\partial t}
    \end{gather*}
    
    \begin{itemize}
        \item Se for exata, seguir o \nameref{par: metodo-resolucao-edo-exata}.
    \end{itemize}
    
    \item Se não for exata, aplicar o método do fator integrante.
    \begin{gather*}
    \mu(t,y) M(t,y) + \mu(t,y) N(t,y) \frac{dy}{dt} = 0
    \end{gather*}

    \item Procurar um fator integrante $\mu(t)$ ou $\mu(y)$.

    \begin{gather*}
    \frac{\frac{\partial M}{\partial y} - \frac{\partial N}{\partial t}}{N} \text{ depende apenas de } t \text{ então existe } \mu(t) \text{ tal que} \\
    \mu'(t) = \mu(t) \frac{\frac{\partial M}{\partial y} - \frac{\partial N}{\partial t}}{N} \\
    \mu'(t) + \left( - \frac{\frac{\partial M}{\partial y} - \frac{\partial N}{\partial t}}{N} \right) \mu(t) = 0 \\
    \implies \mu(t) = e^{\int \frac{\frac{\partial M}{\partial y} - \frac{\partial N}{\partial t}}{N} dt}
    \end{gather*}

    \begin{gather*}
    \frac{\frac{\partial N}{\partial t} - \frac{\partial M}{\partial y}}{M} \text{ depende apenas de } y \text{ então existe } \mu(y) \text{ tal que} \\
    \mu'(y) = \mu(y) \frac{\frac{\partial N}{\partial t} - \frac{\partial M}{\partial y}}{M} \\
    \mu'(y) + \left( - \frac{\frac{\partial N}{\partial t} - \frac{\partial M}{\partial y}}{M} \right) \mu(y) = 0 \\
    \implies \mu(y) = e^{\int \frac{\frac{\partial N}{\partial t} - \frac{\partial M}{\partial y}}{M} dy}
    \end{gather*}

    \medskip

    \item Se existir um fator integrante $\mu$, encontrar a função $\phi(t,y)$, tal que
    \begin{gather*}
    \begin{dcases}
    \frac{\partial \phi}{\partial t} = \mu M \implies \phi = \int \mu M dt + c_1(y) \\
    \frac{\partial \phi}{\partial y} = \mu N \implies \phi = \int \mu N dy + c_2(t)
    \end{dcases}
    \end{gather*}

    Basta calcular uma das equações, e depois determinar o seu termo $c$.

    \item A solução é dada implicitamente por:
    \begin{gather*}
    \phi(t,y) = c, \quad c \in \mathbb{R}
    \end{gather*}

\end{enumerate}




\end{document}

